%% =========================================================================
%% Function field analog --- merged from former companion paper
%% =========================================================================

%% Local macros (kept here so this file is self-contained)
\providecommand{\Fp}{\FF_p}
\providecommand{\Fpt}{\FF_p[t]}
\providecommand{\GF}[2]{\FF_{#1^{#2}}}
\providecommand{\ffseq}{\mathrm{ffSeq}}
\providecommand{\ffprod}{\mathrm{ffProd}}
\providecommand{\irredcount}{\pi_p}
\providecommand{\FFMC}{\mathrm{FF\text{-}MC}}
\providecommand{\FFDH}{\mathrm{FF\text{-}DH}}
\providecommand{\FCD}{\mathrm{FCD}}
\providecommand{\SPV}{\mathrm{SPV}}
\providecommand{\PSCDff}{\mathrm{PSCD}}
\providecommand{\CME}{\mathrm{CME}}
\providecommand{\CCSB}{\mathrm{CCSB}}
\providecommand{\DSL}{\mathrm{DSL}}

\subsection{The Function Field Analog}
\label{sec:function-field}

\begin{roadmap}
Replacing $\ZZ$ by $\Fpt$ and primes by monic irreducible polynomials
gives a clean test bed where \emph{the analytic number theory becomes
free}.  The Weil bound---an open hypothesis over~$\ZZ$---is a theorem of
Weil~\cite{Weil1948} over~$\Fpt$.  Yet the function-field Mullin
Conjecture is not uniformly true: over $\FF_5[t]$ it is \emph{false},
because the greedy sequence from~$t$ is a perpetually irreducible
Sylvester tower (Theorem~\ref{thm:ffmc-false-five}, machine-checked),
and where it is not refuted the residual difficulty is identical to the
integer case: the orbit-specificity barrier.  This subsection develops
the FF-EM sequence, proves the refutation over~$\FF_5$, an unconditional
six-step sieve chain and the $\Phi_3$ exclusion criterion, and presents
three formally verified dead ends that delineate the boundary of current
techniques.  The function-field module of the formalization comprises
19~files and ${\sim}9{,}200$~lines with zero \texttt{sorry}.

\medskip\noindent\textbf{What is and is not machine-checked here (audit
2026-08-17).}\enspace The Weil bound is not formalized: \code{WeilBound},
\code{FFPopulationEquidist}, \code{FFDirichletEquidist},
\code{PopulationMultCCSB}, \code{FFMultiplierCCSB}, \code{FFDSL},
\code{FFCME}, \code{SelectionBiasNeutral} and a few relatives are
\emph{placeholder propositions}---their Lean bodies are \code{True}, they
record intended statements, and every ``theorem'' whose only content is
an implication between them (\code{weil\_implies\_population\_equidist},
\code{weil\_implies\_population\_mult\_ccsb},
\code{selection\_bias\_chain}, clause~(1) of
\code{barrier\_setting\_independent}, \dots) is trivial.  Such items are
marked \placeholder\ below.  The machine-checked content of this
section is: the FF-EM sequence and its structural theory (coprimality
cascade, injectivity, degree growth), the refutation of FF-MC over
$\FF_5[t]$ (\code{EM/FunctionField/StableTower.lean}),
FFDH~$\Rightarrow$~FF-MC, the $\Phi_3$ exclusion, the exact degree telescope
(\S\ref{sec:ff-telescope}), the stochastic FF-MC and its phase
transition, and the density chain whose one
hypothesis \code{FFDirichletDensity} is honestly stated (\code{EM/FunctionField/DensityMC.lean}).
\end{roadmap}

\paragraph{Setting and definitions.}

Let $p$ be a prime.  The ring $\Fpt$ is a Euclidean domain, hence a UFD;
\emph{monic irreducible polynomials} play the role of prime numbers.  We
write $\irredcount(d)$ for the number of monic irreducibles of degree~$d$
over~$\Fp$.  For~$p=2$ the first values are
$\irredcount(1) = 2$ (namely $t$ and $t{+}1$),
$\irredcount(2) = 1$ (namely $t^2{+}t{+}1$), $\irredcount(3) = 2$,
$\irredcount(4) = 3$.

\begin{definition}[FF-EM Sequence ---
\lean{EM/FunctionField/Analog.lean}{FFEMData}]
\label{def:ffem}
Fix a prime~$p$.  An \defname{FF-EM Data} consists of:
$\ffseq(0) = t$ (the smallest monic irreducible);
$\ffprod(0) = t$;
$\ffseq(n{+}1)$ a monic irreducible factor of $\ffprod(n)+1$ of
\emph{minimal degree};
$\ffprod(n{+}1) = \ffprod(n)\cdot\ffseq(n{+}1)$.
\end{definition}

\begin{conjecture}[Function-Field Mullin Conjecture ---
\lean{EM/FunctionField/Analog.lean}{FFMullinConjecture}]
\label{conj:ffmc}
For every monic irreducible $Q \in \Fpt$, there exists $n$ with
$\ffseq(n) = Q$.
\end{conjecture}

Stated for every prime~$p$, the conjecture is \emph{false}.

\begin{theorem}[{FF-MC fails over $\FF_5[t]$ ---
\lean{EM/FunctionField/StableTower.lean}{not\_ffMullinConjecture\_five}}]
\label{thm:ffmc-false-five}
Over $\FF_5[t]$ every Euclid polynomial $\ffprod(n)+1$ of the sequence
seeded at~$t$ is irreducible
(\lean{EM/FunctionField/StableTower.lean}{tower\_euclid\_irreducible}).
Consequently the sequence is
\[
  t,\; t+1,\; t^2+t+1,\; t^4+2t^3+2t^2+t+1,\; g_3(t),\; g_4(t),\;\dots,
  \qquad g_n := \Phi_6^{\,n} + 1,\ \ \Phi_6(y) = y(y+1),
\]
no tie ever occurs, the selected irreducibles are $t$, $t+1$ and exactly
one of each degree $2^k$ for $k \geq 1$, and every other monic
irreducible---every one of degree not a power of two, and the linear
polynomials $t+2$, $t+3$---never appears.  In particular $\neg\FFMC$ for
$p = 5$, for every choice function.
\end{theorem}

\begin{proof}[Proof sketch]
On the autonomous branch $\ffprod(n) = \Phi_6^{\,n}(t)$ and
$\ffprod(n)+1 = g_n(t)$, so the claim is that every level polynomial
$g_n$ is irreducible over~$\FF_5$
(\lean{EM/FunctionField/StableTower.lean}{g\_irreducible}).  In
$L = \overline{\FF}_5$ put $\alpha_0 = -1$ and
$\alpha_{k+1} = 3s_{k+1}+2$ with $s_{k+1}^2 = 1+4\alpha_k$, so that
$\alpha_{k+1}^2+\alpha_{k+1} = \alpha_k$ and $g_n(\alpha_n) = 0$.  With
$q_e = 5^{2^e}$ and $m_e = (q_e-1)/2$ a single induction, using only
that $x \mapsto x^{q_e}$ is a ring endomorphism fixing $\FF_5$ and the
identity $m_{e+1} = m_e(q_e+1)$, gives
$\alpha_k^{q_k} = \alpha_k$, $(1+4\alpha_k)^{m_k} = -1$ and
$(2+4\alpha_k)^{m_k} = -1$
(\lean{EM/FunctionField/StableTower.lean}{tower\_invariant}).  Hence
$s_{k+1}^{q_k} = -s_{k+1}$ and $\alpha_n^{q_{n-1}} = -3s_n+2 \neq \alpha_n$:
the Frobenius orbit of $\alpha_n$ has exactly $2^n$ points
(\lean{EM/FunctionField/StableTower.lean}{minimalPeriod\_alpha}), all
roots of the minimal polynomial of~$\alpha_n$, which therefore has
degree $\geq 2^n = \deg g_n$ and divides~$g_n$; so $g_n$ is the minimal
polynomial.  No Capelli lemma and no norm computation is needed.  The
linear polynomial $t+3$ then witnesses the failure: the only linear
terms are $t$ and $t+1$.
\end{proof}

The mechanism is a \emph{stable tower}.  By Capelli's lemma and the
finite-field square criterion, $g_n$ is irreducible over~$\Fp$ iff
$g_{n-1}$ is and $\Phi_6^{\,n-1}(-1/4)+1$ is a non-square mod~$p$; the
first values of these constants over~$\mathbb{Q}$ are $13/16$, $217/256$,
$57073/65536$---the level constants of \S\ref{sec:level-tower}, made
rational again by the norm.  Call $p$ \emph{exceptional} if $p \equiv 2
\pmod 3$ and every value on the (finite) orbit of $-1/4$ under
$y \mapsto y^2+y$ in $\Fp$ is one less than a non-square.  For $p = 5$
one has $-1/4 = 1$, the orbit is the two-cycle $1,2,1,2,\dots$, and the
shifted values $3, 2$ are both non-squares: $5$ is exceptional.  The
primes $11, 17, 23, 29, 41, 47$ are not (for $17, 23, 29$ already $13$
is a square; $11$ and $41$ fail at level~$4$, $47$ at level~$8$), and
heuristically the exceptional primes are finite in number.  Exactly for
exceptional~$p$ is the seed-$t$ sequence a perpetually irreducible
tower, and then FF-MC fails as in Theorem~\ref{thm:ffmc-false-five}.
The honest function-field conjecture therefore excludes them:

\begin{conjecture}[Function-Field Mullin Conjecture, corrected]
\label{conj:ffmc-corrected}
For every non-exceptional prime~$p$ and every FF-EM data over~$\Fpt$,
every monic irreducible $Q \in \Fpt$ appears.
\end{conjecture}

Two consequences are worth stating.  First, the composite floor
$(\mathrm{C}\infty)$ and MC fail \emph{together} over $\FF_5[t]$: the
perpetually irreducible Sylvester tower that cannot be excluded
over~$\ZZ$ (\S\ref{sec:cinfty-verdict}) is realised, and the growth
constant of \S\ref{sec:ff-telescope} equals~$1$ there
($\deg\ffprod(N) = 2^N$ exactly).  Second, \code{FFPerpetualIrreducibility}
(\code{EM/FunctionField/DegreeTelescope.lean}), treated below as a
hypothesis, is a \emph{theorem} for $p = 5$ from stage~$0$.  In the other
direction, over $\FF_2[t]$ the take-all map $P \mapsto P^2+P$ is additive
and the identity $y^8+y^4+y^2+y+1 = (y^4+y^3+1)(y^4+y^3+y^2+y+1)$ shows
that no four consecutive Euclid polynomials are irreducible, so
$(\mathrm{C}\infty)$ holds there for every seed; and for $p \equiv 1
\pmod 3$ already $\Phi_3$ splits, so no two consecutive Euclid
polynomials are irreducible.  Both floors are machine-checked for every
FF-EM data (any choice function):
\lean{EM/FunctionField/CompositeFloors.lean}{euclid\_three\_reducible\_of\_two}
and
\lean{EM/FunctionField/CompositeFloors.lean}{ffGrowthConstant\_eq\_zero\_of\_two}
for $p = 2$,
\lean{EM/FunctionField/CompositeFloors.lean}{euclid\_succ\_reducible\_of\_one\_mod\_three}
and
\lean{EM/FunctionField/CompositeFloors.lean}{ffGrowthConstant\_eq\_zero\_of\_one\_mod\_three}
for $p \equiv 1 \pmod 3$.  Third, the stable tower over $\FF_5$ is not a
feature of the seed~$t$ alone: for the quadratic seeds $t^2+1$ and
$t^2+2$ every $g_n(t^2+c)$ is irreducible over $\FF_5$
(\lean{EM/FunctionField/QuadraticSeeds.lean}{g\_comp\_irreducible}), so
every admissible factor selection from these seeds follows the tower and
never selects a linear irreducible
(\lean{EM/FunctionField/QuadraticSeeds.lean}{quad\_seed\_perpetual},
\lean{EM/FunctionField/QuadraticSeeds.lean}{quad\_seed\_sel\_natDegree}).
The criterion is the norm criterion again: the seed $t^2+c$ is perpetual
iff the $\Phi_6$-orbit of~$c$ in $\FF_5$ stays inside $\{1,2\}$, and
$c = 0, 3, 4$ break at the first step.  Since a quadratic seed is affinely
conjugate to some $t^2+c$, ten of the twenty-five monic quadratics over
$\FF_5$ are perpetual seeds; for degree $\geq 3$ the Capelli condition is
a value-set condition, not a norm, and is expected to fail within a few
levels.  The failing seeds over $\FF_5$ are the small structured ones.
The five linear seeds fail as well: for every $a \in \FF_5$ each
$g_n(t+a)$ is irreducible
(\lean{EM/FunctionField/LinearSeeds.lean}{g\_comp\_linear\_irreducible},
\lean{EM/FunctionField/LinearSeeds.lean}{lin\_seed\_perpetual}), because
subtracting a constant does not change a Frobenius period.

\paragraph{The small characteristics in full.}
The Frobenius-orbit method is stated once for every prime in
\code{EM/FunctionField/FrobeniusOrbit.lean}: for $\beta \in \overline{\FF}_p$
algebraic, the period of $\beta$ under $x \mapsto x^p$ is at most the degree
of its minimal polynomial
(\lean{EM/FunctionField/FrobeniusOrbit.lean}{minimalPeriod\_le\_natDegree\_minpoly}),
so a monic polynomial with a root whose period equals its degree is
irreducible
(\lean{EM/FunctionField/FrobeniusOrbit.lean}{irreducible\_of\_natDegree\_eq\_minimalPeriod});
conversely a root of an irreducible of degree~$d$ lies in $\FF_{p^d}$
(\lean{EM/FunctionField/FrobeniusOrbit.lean}{pow\_p\_pow\_natDegree\_eq\_self}).
With it, the three smallest characteristics are worked out for every
FF-EM data, i.e.\ for every choice function:
\begin{itemize}[nosep,leftmargin=1em]
\item $\FF_2$.  The polynomials $t^2+t+1$ and $t^4+t+1$ are irreducible
  (roots of period $2$ and $4$), so the first four terms are forced,
  $t,\ t+1,\ t^2+t+1,\ t^4+t+1$, and the fourth Euclid polynomial is the
  tie $(t^4+t^3+1)(t^4+t^3+t^2+t+1)$
  (\lean{EM/FunctionField/CharTwo.lean}{ff\_two\_first\_terms},
  \lean{EM/FunctionField/CharTwo.lean}{euclid\_three\_factor}).  In
  particular the constant~$3$ of the $\FF_2$ floor is attained
  (\lean{EM/FunctionField/CharTwo.lean}{ff\_two\_attains\_three}).
\item $\FF_3$.  Here $\Phi_3 = (y-1)^2$, so an irreducible Euclid polynomial
  is followed by a \emph{perfect square}, $E_{n+1} = (P_n-1)^2$
  (\lean{EM/FunctionField/CharThree.lean}{euclid\_succ\_eq\_sq}); the floor
  holds with constant~$1$
  (\lean{EM/FunctionField/CharThree.lean}{ffGrowthConstant\_eq\_zero}), the
  selected factor after an irreducible stage divides $P_n - 1$, and the
  first five terms are forced: $t,\ t+1,\ t+2,\ t^3+2t+1,\ t^3+2t+2$
  (\lean{EM/FunctionField/CharThree.lean}{ff\_three\_first\_terms}), with
  $\ffprod(2) = t^3 - t$ the product of all linear polynomials.
\item Degrees after autonomous steps.  For $p \equiv 2 \pmod 3$, every
  irreducible factor of the Euclid polynomial following an irreducible one
  has even degree
  (\lean{EM/FunctionField/FrobeniusOrbit.lean}{even\_natDegree\_of\_dvd\_phi3},
  \lean{EM/FunctionField/AutonomousDegrees.lean}{ffSeq\_natDegree\_even\_of\_irreducible}),
  the full form of the $\Phi_3$ exclusion below, unconditionally and for a
  single step; over $\FF_2$, after two irreducible stages every factor has
  degree divisible by~$4$
  (\lean{EM/FunctionField/AutonomousDegrees.lean}{four\_dvd\_ffSeq\_natDegree\_of\_two\_irreducible}).
\end{itemize}

\paragraph{The generic sequence, and a corollary over $\mathbb{Q}$.}
A monic integer polynomial whose reduction mod a prime is irreducible is
irreducible over~$\ZZ$, hence over~$\mathbb{Q}$ by Gauss's lemma.
Theorem~\ref{thm:ffmc-false-five} therefore gives, at no extra cost:

\begin{theorem}[Level polynomials are irreducible over $\mathbb{Q}$ ---
\lean{EM/FunctionField/GenericTower.lean}{gQ\_irreducible},
\lean{EM/FunctionField/GenericTower.lean}{gZ\_irreducible}]
\label{thm:level-irreducible-Q}
For every $n$, the level polynomial $g_n = \Phi_6^{\,n} + 1 \in \ZZ[y]$ is
irreducible over $\ZZ$ and over~$\mathbb{Q}$.  Hence the Galois group acts
transitively on every level of the backward-orbit tree of $-1$ under
$y \mapsto y^2+y$ (\S\ref{sec:level-tower}).
\end{theorem}

This has a reading that unifies the three settings of this paper.  Over
$R[x]$ with the seed~$x$---the \emph{generic point}---the accumulator on
the autonomous branch is $\Phi_6^{\,n}(x)$ and the Euclid polynomial is
$g_n(x)$ (\lean{EM/FunctionField/GenericTower.lean}{iterR}); over
$\mathbb{Q}[x]$ every one of them is irreducible, so \emph{the generic
Euclid--Mullin sequence is a perpetually irreducible Sylvester tower},
$x,\ x+1,\ x^2+x+1,\ x^4+2x^3+2x^2+x+1,\ \dots$, and misses every
irreducible polynomial not of the form~$g_n$.  Every function-field
seed-$t$ sequence is the \emph{reduction mod~$p$} of this generic
sequence, and the integer Euclid--Mullin sequence is its
\emph{specialization} at $x = 2$
($g_0(2) = 3$, $g_1(2) = 7$, $g_2(2) = 43$).  The reduction leaves the
tower at the first level where $g_k$ factors mod~$p$---a Galois-theoretic
event, and never for an exceptional prime such as~$5$.  The specialization
leaves it at stage~$3$, because $g_3(2) = 1807 = 13\cdot 139$ is composite
(\lean{EM/FunctionField/GenericTower.lean}{gZ\_three\_eval\_two})---an
arithmetic event about a \emph{value} of an irreducible polynomial.
Mullin's conjecture over~$\ZZ$ lives entirely in this third regime:
$(\mathrm{C}\infty)$ says that the specialization leaves every tower it
enters, whereas the generic sequence never does.

Two structural differences from the integer setting are immediate.
First, monic irreducibles of the same degree have no canonical total
order, so the ``minimal-degree'' rule may select among several
\emph{competitors}---the analog of the integer $\minFac$ rule loses
uniqueness.  Second, the unit group
$(\Fpt/(Q))^\times \cong \GF{p}{\deg Q}^\times$ is cyclic of order
$p^{\deg Q}-1$, providing richer group-theoretic structure than
$(\ZZ/q\ZZ)^\times$.

\paragraph{Why study the function-field case?}
The Weil bound for character sums over $\Fp$-curves is a
\emph{theorem}~\cite{Weil1948,Deligne1980}, not a conjecture.  This
makes population equidistribution of monic irreducibles across residue
classes unconditional, in contrast to the integer setting where the
analogous result requires the Weighted PNT in arithmetic progressions.
By isolating the \emph{orbit-specificity barrier}---the gap between
``population statistics'' and ``the specific walk
$\ffprod(n) \bmod Q$ visits the residue class $-1$''---we will show
that, for the primes not refuted by Theorem~\ref{thm:ffmc-false-five},
this barrier is the sole remaining obstacle to FF-MC, and that it is
structurally identical to the obstacle for MC over~$\ZZ$.

\paragraph{Basic structural facts.}
The following hold unconditionally for any FF-EM data:

\begin{lemma}[Coprimality cascade ---
\lean{EM/FunctionField/Analog.lean}{coprimality\_cascade}]
\label{lem:ff-coprime}
No monic irreducible $Q$ divides both $\ffprod(n)$ and $\ffprod(n)+1$.
\end{lemma}

\begin{proof}
If $Q$ divides both, then $Q \mid 1$, contradicting $\deg Q \geq 1$.
\end{proof}

\begin{lemma}[Sequence injectivity ---
\lean{EM/FunctionField/Bootstrap.lean}{ffSeq\_injective\_proved}]
\label{lem:ff-injective}
The map $n \mapsto \ffseq(n)$ is injective: no irreducible appears
twice.
\end{lemma}

\begin{lemma}[Degree growth ---
\lean{EM/FunctionField/Analog.lean}{ffProd\_natDegree\_strict\_mono}]
\label{lem:ff-degree-growth}
$\deg(\ffprod(n+1)) > \deg(\ffprod(n))$ for all~$n$.
\end{lemma}

\begin{lemma}[Pool degree to infinity ---
\lean{EM/FunctionField/Finiteness.lean}{ffSeq\_degree\_tendsto\_atTop\_unconditional}]
\label{lem:ff-pool}
$\deg(\ffprod(n)+1) \to \infty$ as $n \to \infty$.
\end{lemma}

\paragraph{The necklace identity.}

The combinatorial heart of the unconditional ANT chain over $\Fpt$ is
the following exact identity.

\begin{theorem}[Necklace identity ---
\lean{EM/FunctionField/NecklaceFormula.lean}{necklace\_identity\_proved}]
\label{thm:necklace}
For every $n \geq 1$,
\[
  \sum_{d \mid n}\, d\cdot\irredcount(d) \;=\; p^n.
\]
\end{theorem}

The name comes from the combinatorial interpretation: $\irredcount(d)$
counts aperiodic necklaces of length~$d$.  The identity is classical
(Moreau, 1872), and is equivalent to the factorization
$t^{p^n}-t = \prod_{\deg Q \mid n} Q$ over $\Fpt$.

\begin{proof}[Proof sketch]
Both sides count the $p^n$ elements of $\GF{p}{n}$, viewed as roots of
$t^{p^n}-t$.  The forward inclusion uses
$Q \mid t^{p^d}-t \mid t^{p^n}-t$ when $\deg Q \mid n$; the reverse
uses the tower law on $[\Fp(\alpha):\Fp]$ to bound the degree of any
minimal polynomial.  The Lean proof spans ${\sim}380$~lines and
exercises Mathlib's Galois-field infrastructure
(\code{GaloisField.card}, \code{minpoly.irreducible},
\code{IntermediateField.finrank\_dvd\_of\_le\_right},
\code{FiniteField.X\_pow\_card\_pow\_sub\_X\_ne\_zero}).
\end{proof}

\begin{corollary}[Density bounds ---
\lean{EM/FunctionField/NecklaceFormula.lean}{necklace\_count\_mul\_le}]
\label{cor:irred-bounds}
For every $n \geq 1$, $\;n\cdot\irredcount(n) \leq p^n$ and
$\irredcount(n) \leq p^n/n$.
\end{corollary}

A finer Möbius inversion gives
$\irredcount(n) = \frac{1}{n}\sum_{d \mid n}\mu(n/d)p^d
  \sim p^n/n + O(p^{n/2}/n)$;
the crude bound suffices for our sieve chain.

\paragraph{The unconditional sieve chain.}

The integer sieve approach to Mullin's Conjecture proceeds via
\[
  \text{Char.\ Cancel.}
  \xRightarrow{\text{GRH}} \text{PNT-AP}
  \xRightarrow{\text{S--W}} \FCD
  \Rightarrow \SPV \Rightarrow \PSCDff
  \Rightarrow \text{Almost-All MC}.
\]
Over $\Fpt$, \emph{every arrow is unconditional}.

\smallskip
\begin{center}
\renewcommand{\arraystretch}{1.2}
\begin{tabular}{lll}
\toprule
\textbf{Step} & \textbf{Over $\ZZ$} & \textbf{Over $\Fpt$} \\
\midrule
Prime counting & PNT & Exact (necklace identity) \\
Primes in APs & Siegel--Walfisz & Existence only ($\irredcount(d)\geq 1$); Weil = \placeholder \\
$\FCD$ & WPNT-in-APs & Linear supply \\
$\SPV$ & FCD + sparsity & Necklace bound \\
$\PSCDff$ & FCD + SPV & Density bound \\
Almost-all MC & PSCD + pigeonhole & PNT + necklace \\
\bottomrule
\end{tabular}
\end{center}

\paragraph{Step 1 --- Counting irreducibles.}
\begin{theorem}[FF irreducible count ---
\lean{EM/FunctionField/FFCharacterSums.lean}{ff\_char\_sum\_cancellation\_proved}]
The necklace identity (Theorem~\ref{thm:necklace}) holds unconditionally
and gives the exact count of monic irreducibles of each degree.
\end{theorem}

To be precise about the analogy: the necklace identity is the
function-field \emph{prime number theorem} (an exact count of primes by
degree), not the analogue of the Riemann Hypothesis for Dirichlet
$L$-functions; that analogue is the Weil bound for $\Fpt$-Dirichlet
characters, which in this development is a \placeholder\ (see the audit note at the start of this section).
What is exact here is counting, not character cancellation.

\paragraph{Step 2 --- PNT in arithmetic progressions.}
\begin{theorem}[FF PNT-in-APs ---
\lean{EM/FunctionField/NecklaceFormula.lean}{ffIrredCount\_pos}]
For every $d \geq 1$, $\irredcount(d) \geq 1$.
\end{theorem}

The proof uses the Galois-field construction: the minimal polynomial of
a primitive element of $\GF{p}{d}$ is monic irreducible of degree~$d$.
The reduction to character cancellation is automatic
(\lean{EM/FunctionField/FFCharacterSums.lean}{ff\_char\_cancel\_implies\_pnt}).

\paragraph{Step 3 --- Forbidden class divergence.}
\begin{property}[FF-FCD ---
\lean{EM/FunctionField/FFCharacterSums.lean}{FFFCD}]
$\sum_{d=1}^{D}\irredcount(d) \geq D$ for all $D \geq 1$.
\end{property}

\begin{theorem}[FF-FCD is unconditional ---
\lean{EM/FunctionField/FFCharacterSums.lean}{ff\_fcd\_proved}]
Each summand is $\geq 1$ by Step~2; the bound is immediate
(\lean{EM/FunctionField/FFCharacterSums.lean}{ff\_pnt\_implies\_fcd}).
\end{theorem}

\paragraph{Step 4 --- Sieve product vanishing.}
\begin{property}[FF-SPV ---
\lean{EM/FunctionField/FFSieve.lean}{FFSieveProductVanishing}]
$n\cdot\irredcount(n) \leq p^n$ for all $n \geq 1$.
\end{property}

\begin{theorem}[FF-SPV is unconditional ---
\lean{EM/FunctionField/FFSieve.lean}{ff\_spv\_proved}]
Immediate from Theorem~\ref{thm:necklace}: the term $n\cdot\irredcount(n)$
is one summand and the total equals $p^n$.  Reduction to FCD:
\lean{EM/FunctionField/FFSieve.lean}{ff\_fcd\_implies\_spv}.
\end{theorem}

\paragraph{Step 5 --- Population sieve confinement decay.}
\begin{property}[FF-PSCD ---
\lean{EM/FunctionField/FFSieve.lean}{FFPSCD}]
$\irredcount(n) \leq p^n/n$ for all $n \geq 1$.
\end{property}

\begin{theorem}[FF-PSCD is unconditional ---
\lean{EM/FunctionField/FFSieve.lean}{ff\_pscd\_proved}]
Divide both sides of Step~4 by~$n$.  Reduction:
\lean{EM/FunctionField/FFSieve.lean}{ff\_spv\_implies\_pscd}.
\end{theorem}

\noindent
As a standalone Prop, \textsf{FFPSCD} is a \emph{counting proxy}: it
records the irreducible-count bound that the sieve consumes, not
confinement decay itself.  The genuine FF confinement decay---for
every proper residue subset, the density of confined starting points
tends to zero---is
\lean{EM/FunctionField/DensityMC.lean}{ff\_density\_pscd}, proved
conditionally on Kornblum's theorem (see the genuine density statement
below).

\paragraph{Step 6 --- Almost-all generalized mixed MC.}
A \emph{mixed} selection at each step picks \emph{any} irreducible factor
of $\ffprod(n)+1$ (not necessarily minimal-degree); the
\emph{generalized} variant starts from an arbitrary monic squarefree
polynomial.  The Step-6 Prop
\lean{EM/FunctionField/FFSieve.lean}{FFAlmostAllGenMixedMC} is again a
counting proxy: it packages the two counting inputs the sieve consumes
(positivity of $\irredcount$ in every degree, and FF-SPV).  The
density-one statement itself is Theorem~\ref{thm:ff-density} below.

\begin{theorem}[FF sieve-chain inputs are unconditional ---
\lean{EM/FunctionField/FFSieve.lean}{ff\_almost\_all\_unconditional}]
\label{thm:ff-almost-all}
For every prime $p$, the conjunction of FF-PNT-in-APs and FF-SPV holds
unconditionally.
\end{theorem}

\begin{keyresult}
\textbf{Full chain witness}
(\lean{EM/FunctionField/FFSieve.lean}{ff\_full\_chain\_unconditional}):
the conjunction
\[
  \text{Char.\ Cancel.} \,\wedge\, \text{PNT-AP} \,\wedge\, \FCD
  \,\wedge\, \SPV \,\wedge\, \PSCDff \,\wedge\,
  \text{Almost-All FF-GenMixedMC}
\]
holds unconditionally for every prime~$p$, where the last conjunct
abbreviates the counting proxy of Step~6.  Zero open hypotheses.
\end{keyresult}

\paragraph{The genuine density statement.}
The density-one theorem that the Step-5/6 proxies stand in for is
formalized in \texttt{EM/FunctionField/DensityMC.lean}, transplanting
the integer-side chain of \S\ref{sec:almost-all-mixed} to $\FF_p[t]$
with one striking simplification: over $\FF_p[t]$ the sieve counting is
\emph{exact}.  Monic polynomials of degree $d \geq \deg N$ distribute
exactly $p^{\,d - \deg N}$ per residue class mod any monic~$N$
(\lean{EM/FunctionField/DensityMC.lean}{ffMonicDeg\_residue\_card},
a division-algorithm bijection with no error term), the congruence
sieve bound
$\leq \bigl(\prod_{P \in S}(1 - p^{-\deg P})\bigr)\,p^{\,d}$
is exact
(\lean{EM/FunctionField/DensityMC.lean}{ffSievedDegCount\_le\_real}),
and at least a quarter of the monics of each degree are squarefree
(\lean{EM/FunctionField/DensityMC.lean}{ffSqfreeDegCount\_quarter},
unconditional).  Trapped starting points are confined to a proper
residue subset mod the target, and a pigeonhole over the finitely many
proper subsets
(\lean{EM/FunctionField/DensityMC.lean}{ff\_trapped\_le\_sum\_confined})
reduces the density statement to per-subset confinement decay.  The
sole analytic input is the FF analogue of Dirichlet's theorem in the
divergence form:

\begin{property}[FF Dirichlet density (Kornblum~\cite{Kornblum1919}) ---
\lean{EM/FunctionField/DensityMC.lean}{FFDirichletDensity}]
For every monic irreducible $Q$ and every $a$ with $Q \nmid a$,
\[
  \sum_{\substack{P \text{ monic irreducible},\ P \equiv a \ (\mathrm{mod}\ Q) \\ \deg P \leq Y}}
    p^{-\deg P} \;\longrightarrow\; \infty
  \qquad (Y \to \infty).
\]
\end{property}

\begin{theorem}[Almost-all FF-GenMixedMC, genuine density ---
\lean{EM/FunctionField/DensityMC.lean}{ff\_almost\_all\_genmixed\_density}]
\label{thm:ff-density}
Assume \textup{FFDirichletDensity}.  Then for every monic irreducible
target~$Q$: among monic squarefree $m$ with $1 \leq \deg m \leq n$, the
proportion of starting points with $Q \nmid m$ from which $Q$ is not
tree-reachable tends to~$0$ as $n \to \infty$.  (The restriction
$Q \nmid m$ is necessary: if $Q \mid m$ then $Q$ divides every mixed
walk product from~$m$, hence never divides $\mathrm{walk}+1$.)
\end{theorem}

\noindent
The contrast with the integer side is instructive: over~$\ZZ$ the
corresponding divergence input ($\sum_{p \equiv a\,(q)} 1/p = \infty$)
is proved unconditionally from the Mathlib $L$-function machinery
(\S\ref{sec:almost-all-mixed}), while Mathlib has no polynomial
Dirichlet $L$-functions, so Kornblum's 1919 theorem---strictly easier
than its integer counterpart, as the relevant $L$-functions are
polynomials in $p^{-s}$---remains an isolated named hypothesis.
Formalizing it is a natural future campaign.

\paragraph{The factor tree and structural prerequisites.}
Independent of the sieve chain, the factor-tree structure (any valid
\emph{mixed} selection) admits the following formally verified
properties:

\begin{theorem}[Factor-tree structure ---
\lean{EM/FunctionField/Master.lean}{theorem\_A\_factor\_tree}]
\label{thm:ff-A}
For any valid mixed selection from any monic squarefree start,
\begin{enumerate}[nosep,label=(\alph*)]
\item the walk explores infinitely many distinct monic irreducibles
  (\lean{EM/FunctionField/StochasticMC.lean}{ff\_mixed\_explores\_infinitely\_many});
\item the factor pool degree $\deg(\mathrm{acc}(n)+1) \to \infty$
  (\lean{EM/FunctionField/PopulationEquidist.lean}{ff\_factor\_pool\_degree\_grows});
\item selection is injective
  (\lean{EM/FunctionField/FactorTree.lean}{ffMixedSel\_injective});
\item walk products are injective
  (\lean{EM/FunctionField/FactorTree.lean}{ffMixedWalkProd\_injective});
\item the start is never recapturable
  (\lean{EM/FunctionField/FactorTree.lean}{start\_not\_capturable}).
\end{enumerate}
\end{theorem}

\paragraph{Mixed walks: what is and is not proved.}

\begin{theorem}[Degree growth along every mixed walk ---
\lean{EM/FunctionField/StochasticMC.lean}{stochastic\_mc\_unconditional}]
\label{thm:ff-stoch}
For any monic irreducible target~$Q$, any monic non-constant start, and
any valid mixed walk (any admissible choice of factor at each step),
there is a stage~$n$ at which the Euclid polynomial has degree
$\geq \deg Q$.
\end{theorem}
\noindent This is all that the Lean name \code{FFStochasticMC} asserts:
the factor pool becomes large enough for $Q$ to be \emph{a potential}
factor.  It says nothing about probability or capture; whether the
greedy walk selects~$Q$ is $\FFMC$ itself, and the probabilistic
``with any $\varepsilon > 0$ of non-greedy selection every target is
captured'' is a heuristic that is \textbf{not} formalized.  Earlier
drafts overstated this theorem and its companion
\lean{EM/FunctionField/StochasticMC.lean}{ff\_phase\_transition\_unconditional}
(``phase transition at $\varepsilon = 0$''), whose Lean content is the
pair \emph{(degree grows along every mixed walk, degree grows along the
greedy walk)}; we no longer use the phrase.
\label{thm:ff-phase}

\paragraph{The bootstrap: from dynamical hitting to FF-MC.}

\begin{property}[FF Dynamical Hitting (FFDH) ---
\lean{EM/FunctionField/Bootstrap.lean}{FFDynamicalHitting}]
\label{prop:ffdh}
For every FF-EM data and every monic irreducible $Q$ not in the
sequence, there exist cofinally many indices~$n$ at which
$Q \mid \ffprod(n)+1$.
\end{property}

\begin{theorem}[One-hypothesis reduction ---
\lean{EM/FunctionField/Bootstrap.lean}{ff\_dh\_implies\_ffmc}]
\label{thm:ff-bootstrap}
$\FFDH \;\Longrightarrow\; \FFMC$ (using the unconditional finiteness of
irreducibles per degree,
\lean{EM/FunctionField/Finiteness.lean}{ff\_dh\_implies\_ffmc\_unconditional}).
\end{theorem}

The proof is by strong induction on degree.  At degree~$d_0$, the
hypothesis gives MC below~$d_0$; a uniform sieve gap exhausts all
sub-degree irreducibles past some $N_0$; FFDH provides cofinally many
hits past $N_0$; injectivity makes each hit consume a distinct
\emph{competitor} of the same degree; finiteness per degree then forces
$Q$ itself to be selected (\emph{competitor exhaustion},
\lean{EM/FunctionField/Bootstrap.lean}{ff\_element\_capture}).  This is
the function-field analog of the integer bootstrap
\code{dynamical\_hitting\_implies\_mullin}; the only difference is the
final exhaustion step, which over $\ZZ$ collapses because primes are
totally ordered.  Read contrapositively with
Theorem~\ref{thm:ffmc-false-five}: $\FFDH$ is \emph{false} for $p = 5$;
the walk of the seed-$t$ sequence modulo $t+3$ never returns to~$-1$.

\paragraph{The \texorpdfstring{$\Phi_3$}{Phi3} exclusion criterion.}

We now develop a structural obstruction with no integer analog.  When
all irreducible factors of $\ffprod(n)+1$ have degree~$1$, the walk
$\ffprod(n) \bmod Q$ for a degree-one target $Q = (t-a)$ follows the
\emph{autonomous map} $f(w) = w(w+1)$ on $\Fp$.  The target $Q$ divides
$\ffprod(n)+1$ iff the walk position $w_n = -1$.  Hence degree-one
capture reduces to: \emph{does the orbit of $f$ reach~$-1$?}

\begin{definition}[Autonomous map ---
\lean{EM/FunctionField/AutonomousMap.lean}{ffAutonomousMap}]
\label{def:autonomous}
For a prime $p$, $f \colon \Fp \to \Fp$ is given by $f(w) = w(w+1)$.
The \emph{orbit} from $a \in \Fp$ is $a, f(a), f^2(a), \ldots$
\end{definition}

\begin{lemma}[Absorption at zero ---
\lean{EM/FunctionField/AutonomousMap.lean}{ffAutonomousMap\_neg\_one}]
$f(-1) = 0$ and $f(0) = 0$; once the orbit hits $0$ it stays there
(\lean{EM/FunctionField/AutonomousMap.lean}{ffAutonomousOrbit\_absorbed}).
\end{lemma}

\begin{lemma}[$\Phi_3$ criterion ---
\lean{EM/FunctionField/AutonomousMap.lean}{ffAutonomousMap\_eq\_neg\_one\_iff}]
\label{lem:phi3}
$f(w) = -1 \iff w^2+w+1 = 0$.  The preimage of $-1$ is the root set of
the third cyclotomic polynomial $\Phi_3(w)$.
\end{lemma}

\begin{lemma}[Cube-root condition ---
\lean{EM/FunctionField/AutonomousMap.lean}{pow\_three\_eq\_one\_of\_phi3}]
If $w^2+w+1 = 0$ in $\Fp$, then $w^3 = 1$.
\end{lemma}

\begin{theorem}[$\Phi_3$ has no roots when $p \equiv 2 \pmod 3$ ---
\lean{EM/FunctionField/AutonomousMap.lean}{phi3\_no\_roots}]
\label{thm:phi3-no-roots}
For $p \equiv 2 \pmod 3$ with $p \geq 5$, $\;w^2 + w + 1 \neq 0$ for all
$w \in \Fp$.
\end{theorem}

\begin{proof}
Suppose $w^2+w+1 = 0$ in $\Fp$.  Then $w^3 = 1$, and $w \neq 0,1$ (each
substitution gives a contradiction with $p \geq 5$).  Hence $\mathrm{ord}(w) = 3$
in $\Fp^\times$, so $3 \mid p-1$ by Lagrange.  But $p \equiv 2 \pmod 3$
forces $p - 1 \equiv 1 \pmod 3$, contradiction.
\end{proof}

\begin{theorem}[Unreachability of $-1$ ---
\lean{EM/FunctionField/AutonomousMap.lean}{ff\_neg\_one\_unreachable}]
\label{thm:ff-unreach}
For $p \equiv 2 \pmod 3$ with $p \geq 5$ and any starting value
$a \neq -1$ in $\Fp$, the orbit $f^n(a)$ never equals $-1$.
\end{theorem}

\begin{corollary}[Visited at most once ---
\lean{EM/FunctionField/AutonomousMap.lean}{ff\_neg\_one\_visited\_at\_most\_once}]
For the same primes, $-1$ is visited only at step~$0$ (and only if
$a = -1$); for all $n \geq 1$, the orbit value is never $-1$.
\end{corollary}

\begin{remark}
The criterion fails to apply for $p = 2$ (where $\Phi_3 = t^2+t+1$ is
itself irreducible) and $p = 3$ (where $\Phi_3$ has the double root
$w = 1$).  Primes $p \equiv 2 \pmod 3$ with $p \geq 5$ include
$5, 11, 17, 23, 29, 41, 47, \dots$
\end{remark}

\begin{remark}[The criterion is unconditional for $p = 5$]
The hypothesis of the criterion---that every factor of $\ffprod(n)+1$
has degree one, so that the walk follows the autonomous map---is a
special case of perpetual irreducibility, which for $p = 5$ and the
seed~$t$ is a theorem at every stage (Theorem~\ref{thm:ffmc-false-five}).
There the exclusion says outright that the linear irreducibles $t+2$ and
$t+3$ are never selected: their walks start at $2, 3 \neq -1$ and follow
$w \mapsto w(w+1)$ forever, and $\Phi_3$ has no root in~$\FF_5$.
\end{remark}

\paragraph{Integer non-analog.}
Over $\ZZ$ the multiplier $a_{n+1} \bmod q$ is \emph{not} determined by
$P(n) \bmod q$ alone---it depends on the global factorization of
$P(n)+1$.  No autonomous map exists, so the $\Phi_3$ criterion has no
direct integer analog.

\paragraph{The cyclotomic structure (Dead End \#129).}
Over $\FF_2$, the product of all monic irreducibles of degree dividing
$d=2$ is $t(t+1)(t^2+t+1) = t^4+t$, and adding $1$ gives the irreducible
$t^4 + t + 1$ (a factor of $\Phi_{15}$; the fifth cyclotomic polynomial
$\Phi_5 = t^4+t^3+t^2+t+1$ appears one stage later, as one of the two
tied quartic factors of $\ffprod(3)+1 = t^8+t^4+t^2+t+1$).  Its Galois
group is $\mathrm{Gal}(\FF_{2^4}/\FF_2) \cong \ZZ/4\ZZ$---abelian,
not~$S_4$.
More generally, the classical identity
$\prod_{\deg Q \mid d}Q = t^{p^d}-t$ forces the greedy FF-EM to build
products dividing $t^{p^d}-t$ for some~$d$, so $\ffprod(n)+1$ has
splitting field inside $\GF{p}{N}$ for some~$N$ and the Galois group is
\emph{cyclic}, generated by Frobenius.  Thus the ``large monodromy''
hypothesis is structurally false for the greedy FF-EM, witnessed by
\lean{EM/FunctionField/Analog.lean}{ff\_cyclotomic\_dead\_end}.

\paragraph{The orbit-specificity boundary.}

We now establish that the orbit-specificity barrier is the sole
remaining obstacle to FF-MC, and that it is structurally identical to
the barrier over~$\ZZ$.  The central dichotomy is between two types of
equidistribution:

\begin{itemize}[nosep,leftmargin=1em]
\item \emph{Population equidistribution}: among all monic irreducibles
  of degree~$d$, residues mod~$Q$ are equidistributed across
  $|\GF{p}{\deg Q}^\times|$ classes, with error $O(p^{d/2})$ from the
  Weil bound.  \textbf{Unconditional} over $\Fpt$.
\item \emph{Orbit equidistribution}: the specific walk
  $\ffprod(n) \bmod Q$ visits each residue class with frequency
  $1/(p^{\deg Q}-1)$ (the function-field orbit hypothesis, the analogue
  of CME).  \textbf{Open}.
\end{itemize}

\begin{remark}[Population CCSB from Weil \placeholder]
Mathematically, the Weil bound gives population-level multiplier CCSB
unconditionally, and the intended hierarchy
population~$\to$~fiber~$\to$~selection~$\to$~orbit is recorded in
\code{weil\_implies\_population\_mult\_ccsb} and
\code{selection\_bias\_chain} (\code{EM/FunctionField/MultiplierCCSB.lean}).
These are placeholder propositions: nothing about the Weil bound is
machine-checked, and the two ``theorems'' are implications between
\code{True}s.
\end{remark}

\paragraph{Dead End \#127 --- Weil bound insufficient.}

\begin{theorem}[Weil does not close the gap ---
\lean{EM/FunctionField/OrbitBarrier.lean}{dead\_end\_127\_weil\_not\_sufficient}]
\label{thm:de127}
The Weil bound implies population equidistribution (clause~(1) of the
Lean statement is a placeholder implication, \placeholder) but does not
imply $\FFDH$, because:
(i) Weil controls $\sum_{f\;\mathrm{monic},\;\deg f = d}\chi(f \bmod Q)$,
i.e.\ averages over \emph{all} polynomials of degree~$d$;
(ii) $\FFDH$ requires $\sum_{n<N}\chi(\ffprod(n) \bmod Q)$, a sum along
a \emph{specific deterministic orbit};
(iii) $\ffprod(n)+1$ is the unique element of its congruence class, not
a random draw.
\end{theorem}

\paragraph{Dead End \#129 --- Large monodromy structurally false.}

\begin{property}[FF Large Monodromy (FFLM) ---
\lean{EM/FunctionField/Analog.lean}{FFLargeMonodromy}]
The Galois group of $\ffprod(n)+1$ eventually contains the alternating
group: $|\mathrm{Gal}| \geq (\deg \ffprod(n))!/2$ for all sufficiently
large~$n$.
\end{property}

\begin{theorem}[FFLM is structurally false ---
\lean{EM/FunctionField/OrbitBarrier.lean}{dead\_end\_129\_monodromy\_dead}]
\label{thm:de129}
The greedy FF-EM produces products dividing $t^{p^d}-t$.  Adding $1$
gives a polynomial whose splitting field lies in $\GF{p}{N}$ for some
$N$, hence Galois group $\cong \ZZ/N\ZZ$, generated by Frobenius---
\emph{abelian}, not $A_d$.  The monodromy approach reduces to abelian
Galois theory $=$ Dirichlet characters $=$ Weil $=$ population
equidistribution; the route is circular.
\end{theorem}

\paragraph{Dead End \#130 --- Generation $\neq$ coverage.}
A natural hope is that subgroup escape (multipliers generate the full
group) might imply dynamical hitting (the walk visits every element).
Both PRE (over~$\ZZ$) and Weil (over~$\Fpt$) prove subgroup escape.
The following counterexample shows this is insufficient.

\begin{theorem}[Generation $\neq$ coverage ---
\lean{EM/FunctionField/OrbitBarrier.lean}{dead\_end\_130\_generation\_ne\_coverage}]
\label{thm:de130}
On $\ZZ/4\ZZ$, the steps $\{1, 3\}$ generate the full group, yet the
\emph{alternating} walk $0, 1, 0, 1, \ldots$ visits only $\{0, 1\}$,
missing $\{2, 3\}$.  Witnessed in
\lean{EM/Group/CyclicWalkCoverage.lean}{alternating\_walk\_misses\_two};
the converse direction (cancellation $\Rightarrow$ coverage via Weyl)
is
\lean{EM/Group/CyclicWalkCoverage.lean}{walk\_cancel\_implies\_covers}.
\end{theorem}

\paragraph{The barrier theses.}

\begin{theorem}[Orbit barrier thesis ---
\lean{EM/FunctionField/OrbitBarrier.lean}{orbit\_barrier\_thesis}]
\label{thm:barrier-thesis}
The orbit-specificity barrier has the following formally verified
properties:
\begin{enumerate}[nosep,label=(\roman*)]
\item All structural prerequisites for FF-MC are unconditional:
  injectivity, degree growth, finiteness per degree, coprimality
  cascade.
\item $\FFDH \Rightarrow \FFMC$ (Theorem~\ref{thm:ff-bootstrap}).
\item $\FFDH$ is the sole remaining hypothesis for non-exceptional
  primes; for $p = 5$ both $\FFDH$ and $\FFMC$ are false
  (Theorem~\ref{thm:ffmc-false-five}).
\item The barrier is not closed by population statistics
  (Theorem~\ref{thm:de127}; in Lean this clause is the placeholder
  implication \code{WeilBound}~$\to$~\code{FFPopulationEquidist}
  \placeholder).
\item The barrier is not closed by Galois theory
  (Theorem~\ref{thm:de129}; the surviving formal content is the
  coprimality cascade).
\item The barrier is not closed by algebraic generation
  (Theorem~\ref{thm:de130}).
\item Along every mixed walk the degree grows past any target
  (Theorem~\ref{thm:ff-stoch}); the probabilistic ``capture for
  $\varepsilon > 0$'' is a heuristic, not a theorem.
\item The barrier is sharp at $\varepsilon = 0$
  (Theorem~\ref{thm:ff-phase}).
\end{enumerate}
The final Lean clause, \code{PopulationMultCCSB}~$\to$
\code{SelectionBiasNeutral}~$\to$~\code{FFMultiplierCCSB}, is a
placeholder implication \placeholder.
\end{theorem}

\begin{theorem}[Setting independence ---
\lean{EM/FunctionField/OrbitBarrier.lean}{barrier\_setting\_independent}]
\label{thm:setting-independent}
\begin{enumerate}[nosep,label=(\roman*)]
\item Over $\Fpt$: ANT is free (Weil $\Rightarrow$ equidistribution of
  irreducibles of large degree modulo~$Q$; in the formalization this
  clause is a placeholder implication \placeholder\ and proves nothing).
\item Over $\Fpt$: orbit specificity remains ($\FFDH$ open for
  non-exceptional~$p$, false for $p = 5$; $\FFDH \Rightarrow \FFMC$).
\item Over $\Fpt$: stochastic perturbation dissolves the barrier;
  the deterministic limit does not.
\end{enumerate}
\end{theorem}

The function-field setting eliminates the ``easy'' obstacle (ANT)
and adds new structural information ($\Phi_3$ exclusion), but the
``hard'' obstacle---orbit specificity, i.e.\ CME and its analogue---is
identical to the integer case.  (The population-level $\minFac$
statement PE fails over $\Fpt$ for the same reason as over $\ZZ$: the
class density of the least-degree irreducible factor among $Q$-rough
polynomials is a head-dominated convergent series, Dead End~\#160; what
Weil supplies is the counting statement for irreducibles of a fixed
large degree, which is not PE\@.)  This proves that, wherever the
conjecture is not simply false (Theorem~\ref{thm:ffmc-false-five}), the
difficulty of Mullin's Conjecture lies \emph{entirely} in the
orbit-specificity barrier and not in the analytic number theory.

\paragraph{The degree telescope.}
The growth axis of \S\ref{sec:defect-telescope} transfers to $\Fpt$, and
it is the one place where the function-field model is unambiguously better
behaved: over $\Fpt$ every approximation in the integer telescope
disappears.  Because that comparison belongs with the telescope rather than
with the function-field programme, it is carried out in
\S\ref{sec:ff-telescope}; the conclusion, that the model is \emph{exactly}
as stuck, is recorded in the table below.


\paragraph{Side-by-side summary.}

\begin{center}
\renewcommand{\arraystretch}{1.2}
\begin{tabular}{lcc}
\toprule
\textbf{Aspect} & \textbf{Over $\ZZ$} & \textbf{Over $\Fpt$} \\
\midrule
Prime equidist.\ mod $q$ (ANT) & Mertens form proved & \textbf{Free} (Weil) \\
Subgroup escape (SE) & Proved (PRE) & Proved (Weil) \\
Coprimality cascade & Proved & Proved \\
$\CME \Rightarrow \MC$ chain & Proved & Proved \\
Orbit specificity (CME / FF analogue) & \textbf{Open} & \textbf{Open} \\
Walk character sums & Not standard & Not standard \\
Stochastic MC ($\varepsilon > 0$) & Proved & Proved \\
$\Phi_3$ exclusion & --- & \textbf{Proved} \\
Defect telescope & Approximate & \textbf{Exact} \\
Growth constant $C = 0 \Leftrightarrow (\mathrm{C}\infty)$ & Proved & Proved \\
$(\mathrm{C}\infty)$ itself & \textbf{Open} & \textbf{False} for $p=5$, seed~$t$ (open otherwise) \\
The conjecture itself & \textbf{Open} & \textbf{False} for $p=5$ (Thm.~\ref{thm:ffmc-false-five}) \\
\bottomrule
\end{tabular}
\end{center}

\paragraph{Open problems specific to FF-MC.}

We close with a few open problems particular to the function-field
analog (the integer counterparts are catalogued in~\S\ref{sec:hard}).

\begin{frontierbox}
\textbf{Deterministic FF-MC.}
For $p = 5$ the answer is \emph{no} (Theorem~\ref{thm:ffmc-false-five}),
and the same holds for every exceptional prime.  For non-exceptional~$p$:
does every monic irreducible over $\Fpt$ eventually appear in the greedy
FF-EM sequence?  Equivalently, does $\FFDH$ hold?  There $\FFDH$ is the
sole remaining hypothesis: the stochastic version ($\varepsilon > 0$) is
proved (Theorem~\ref{thm:ff-stoch}), and the open content is exactly the
$\varepsilon = 0$ boundary.  Two further questions: is the set of
exceptional primes finite (heuristically yes; is $5$ the only one?), and
does the seed-$t$ sequence over a non-exceptional prime, once it leaves
the autonomous branch, ever return to it for a long stretch?
\end{frontierbox}

\begin{frontierbox}
\textbf{The $\varepsilon \to 0$ limit.}
Heuristically, $\varepsilon$-perturbed selection captures every target
almost surely for every fixed $\varepsilon > 0$ (not formalized; see
Theorem~\ref{thm:ff-stoch}).  Even granting it, the integer case shows
the limit is not a route: the $\varepsilon$-family is a faithful
reformulation of MC, not a relaxation (Dead End~\#159), and the same
holds over $\FF_p[t]$.
\end{frontierbox}

\begin{frontierbox}
\textbf{Large monodromy beyond cyclotomic starts.}
Dead End~\#129 shows the \emph{greedy} FF-EM produces cyclotomic (hence
abelian) products.  Does the obstruction persist for non-greedy starting
points?  Cohen's theorem~\cite{Cohen1981} shows a random monic
polynomial of degree~$d$ has Galois group $S_d$ with probability
tending to~$1$.  Is the recursive FF-EM structure forcing cyclotomic
products an artifact of the greedy minimal-degree selection?
\end{frontierbox}

\begin{frontierbox}
\textbf{Integer analog of the $\Phi_3$ criterion.}
The $\Phi_3$ exclusion (Theorem~\ref{thm:ff-unreach}) uses the absence
of cube roots of unity in $\Fp$.  Over $\ZZ$ the walk does not follow
an autonomous map---can systematic exclusion criteria be developed
nonetheless for specific $(q,a)$ pairs?  Cox--van der
Poorten~\cite{CoxVdP1968} suggested structural reasons for the
difficulty of the prime~$41$; the function-field $\Phi_3$ result is a
formal analog.
\end{frontierbox}

\begin{frontierbox}
\textbf{Selection bias neutrality}
(\code{SelectionBiasNeutral} in \code{EM/FunctionField/MultiplierCCSB.lean}
is a placeholder proposition; the question is stated here in prose).
Does greedy minimal-degree selection preserve character-sum cancellation?
A positive answer would close the orbit-specificity gap for both FF-MC
and MC simultaneously, since the question concerns the selection
mechanism, not the ambient ring.  This is the cleanest formulation of
the barrier; the integer analog is CME.
\end{frontierbox}

\paragraph{Headline summary.}
The function-field portion of the formalization is collected in the
landscape theorem
\lean{EM/FunctionField/Master.lean}{ff\_unconditional\_grand\_summary},
which assembles theorems~A, B, C, D
(\lean{EM/FunctionField/Master.lean}{theorem\_A\_factor\_tree},
 \lean{EM/FunctionField/Master.lean}{theorem\_B\_stochastic\_mc},
 \lean{EM/FunctionField/Master.lean}{theorem\_C\_phase\_transition},
 \lean{EM/FunctionField/Master.lean}{theorem\_D\_sieve\_chain})
into a single conjunction with zero open hypotheses.  Together they
witness:
\begin{itemize}[nosep,leftmargin=1em]
\item the unconditional six-step sieve chain
  (Theorem~\ref{thm:ff-almost-all});
\item stochastic FF-MC and the $\varepsilon = 0$ phase transition
  (Theorems~\ref{thm:ff-stoch}, \ref{thm:ff-phase});
\item the necklace identity (Theorem~\ref{thm:necklace});
\item the $\Phi_3$ exclusion criterion (Theorem~\ref{thm:ff-unreach});
\item the orbit-barrier and setting-independence theses
  (Theorems~\ref{thm:barrier-thesis}, \ref{thm:setting-independent}).
\end{itemize}
Separately, and not part of that conjunction, the refutation over
$\FF_5[t]$ (Theorem~\ref{thm:ffmc-false-five}) shows that the
function-field conjecture is a statement about non-exceptional primes.
The function-field treatment thus delivers the cleanest available
diagnosis: \emph{where the conjecture is not false outright, the
difficulty of Mullin's Conjecture is the orbit-specificity barrier, not
the analytic number theory}---and where it is false, the failure is the
one branch, a perpetually irreducible Sylvester tower, that no method
can exclude over~$\ZZ$.
